\documentclass[11pt,a4paper]{article}
\usepackage[utf8]{inputenc}
\usepackage[T1]{fontenc}
\usepackage{amsmath,amssymb,booktabs,geometry,hyperref}
\geometry{margin=2.4cm}
\hypersetup{colorlinks=true,linkcolor=black,citecolor=black,urlcolor=blue}

\title{A seedless loop reading of Koide's relation,\\ and dated predictions for the neutrino masses}
\author{Anthony Charretier\\ \small Independent researcher\\ \small \url{https://github.com/innermost47/grm-core}}
\date{4 November 2026}

\begin{document}
\maketitle

\begin{abstract}
Koide's relation $Q=\sum m/(\sum\sqrt m)^2=2/3$ holds for the charged leptons to a few parts in $10^{6}$ and has no
accepted explanation. We read it as a balance with no seed and no scale: the amplitudes $r_n=\sqrt{m_n}$ split into a
shared part (the ``full'') and deviations that sum to zero (the ``void''), and the relation is exactly
$|\text{void}|=|\text{full}|$. Written as a network of relations and solved by weighted least squares, this single
condition reproduces each charged lepton mass from the other two (0.4\,$\sigma$), and fixes the phase of the three
amplitudes to $0.222222047$, which differs from $2/9$ by $1.75\times10^{-7}$. Applied without change to the neutrinos,
with only the measured mass-squared differences as input, the loop closes only if one amplitude is in opposition and
only in the normal ordering. It gives $m_{1,2,3}=0.365,\ 8.685,\ 50.30$\,meV and $\Sigma m_\nu=59.35\pm0.30$\,meV, hence
$m_\beta=8.84$\,meV and $1.1\le m_{\beta\beta}\le4.1$\,meV. These values were archived with a timestamp before the
measurements that can refute them. A second, more speculative bet on the phase of the neutrino amplitudes is given
together with its first check against JUNO.
\end{abstract}

\section{Koide's relation as a balance}
Koide's relation~\cite{koide} concerns three masses. For three positive masses define $r_n=\sqrt{m_n}$ and split the vector $\mathbf r$ into its component along
$(1,1,1)$ and the remainder,
\begin{equation}
  F=3\bar r^{\,2},\qquad V=\sum_n (r_n-\bar r)^2,\qquad \bar r=\tfrac13\sum_n r_n .
\end{equation}
The deviations $r_n-\bar r$ sum to zero. Since $\sum m_n=F+V$ and $(\sum r_n)^2=3F$,
\begin{equation}
  Q=\frac{\sum m_n}{(\sum r_n)^2}=\frac13\left(1+\frac VF\right),
\end{equation}
so $Q=2/3$ is exactly $V=F$: the shared part and the part that cancels out carry the same weight. Equivalently
$\mathbf r$ makes an angle of $45^\circ$ with $(1,1,1)$, as noted by Foot~\cite{foot}. In the parametrization
$r_n=\mu\,\bigl(1+k\sqrt2\cos(\delta+2\pi n/3)\bigr)$ one has $Q=(1+k^2)/3$, so the balance is $k=1$ and leaves the
phase $\delta$ free~\cite{brannen,zenczykowski}.

The reading we propose is a loop with no seed: no mass is a starting point, no scale enters, and the only condition is
that neither part prevails. The condition fixes one ratio among three masses, not all of them.

\section{Charged leptons}
With the PDG~2024 masses~\cite{pdg}, $Q=0.6666645$ and $|Q-2/3|=2.2\times10^{-6}$, compatible with zero
($0.4\,\sigma$, the uncertainty coming almost entirely from $m_\tau$). We solved the loop as a network of relations
(nodes $m_n$ and $r_n$, relations $m_n=r_n^2$ and $V=F$) with the GRM solver~\cite{grm}. Given one mass, the other
two remain undetermined, as they should. Given two, the third follows:
\begin{center}
\begin{tabular}{lll}
\toprule
predicted & from & loop \quad (measured) \\
\midrule
$m_\tau$ & $m_e,\ m_\mu$ & $1776.969$\,MeV \quad ($1776.93\pm0.09$) \\
$m_\mu$ & $m_e,\ m_\tau$ & $105.656\pm0.006$\,MeV \quad ($105.6583755$) \\
$m_e$ & $m_\mu,\ m_\tau$ & $0.51087\pm0.00030$\,MeV \quad ($0.51099895$) \\
\bottomrule
\end{tabular}
\end{center}
Each prediction lies $0.4\,\sigma$ from the measurement. This is Koide's known success; we report it as the check
that the network reproduces it.

\paragraph{The phase.} Imposing $Q=2/3$ exactly and using only $m_e$ and $m_\mu$, the phase is
$\delta=0.222222047$ (cf. Brannen~\cite{brannen}), against $2/9=0.222222222$: a difference of $1.75\times10^{-7}$
(8 parts in $10^7$). The two statements $Q=2/3$ and $\delta=2/9$ therefore cannot both be exact for these masses; they
would predict $m_\tau=1776.969$ and $1776.967$\,MeV, indistinguishable today. Among the 95 fractions $p/q$ with
$q\le12$ in $[0,2\pi/3)$, the expected number within $1.75\times10^{-7}$ of a given value is $1.6\times10^{-5}$.
The value $2/9$ was, however, noticed after looking at these masses, and this estimate does not correct for it.

\section{Neutrinos}
Only the mass-squared differences are measured. With the PDG~2024 values~\cite{pdg}
($\Delta m^2_{21}=7.53\times10^{-5}$\,eV$^2$, $\Delta m^2_{32}=2.455\times10^{-3}$\,eV$^2$ in the normal ordering,
$-2.529\times10^{-3}$ in the inverted one) the three relations
\begin{equation}
  r_2^4-r_1^4=\Delta m^2_{21},\qquad r_3^4-r_2^4=\Delta m^2_{32},\qquad V=F
\end{equation}
were solved from many starting points and for every choice of signs of the amplitudes. With all amplitudes of the same
sign the loop does not close. It closes with one amplitude in opposition to the other two, as proposed for the
neutrinos by Brannen~\cite{brannen}:
\begin{center}
\begin{tabular}{lcccc}
\toprule
ordering & $m_1$ (meV) & $m_2$ (meV) & $m_3$ (meV) & $\Sigma m_\nu$ (meV) \\
\midrule
normal & $0.365$ & $8.685$ & $50.30$ & $59.35\pm0.30$ \\
inverted & $49.62$ & $50.38$ & $2.94$ & $102.9\pm0.6$ \\
\bottomrule
\end{tabular}
\end{center}
The uncertainty follows from the measured differences. The inverted solution exceeds the DESI~2024 bound
$\Sigma m_\nu<72$\,meV (95\,\%, with CMB, $\Lambda$CDM)~\cite{desi}, so the loop predicts the normal ordering, with an
almost massless lightest neutrino. With the mixing angles of PDG~2024 ($\sin^2\theta_{12}=0.307\pm0.013$,
$\sin^2\theta_{13}=0.0219\pm0.0007$), the effective masses are
\begin{equation}
  m_\beta=\Bigl(\sum_i|U_{ei}|^2m_i^2\Bigr)^{1/2}=8.84\ \text{meV},\qquad
  1.1\ \text{meV}\le m_{\beta\beta}\le4.1\ \text{meV},
\end{equation}
the range of $m_{\beta\beta}$ covering the two unknown Majorana phases and the uncertainties of the angles.

\section{A more speculative bet: the phase on a Cantor cascade}
The terms removed by the Cantor construction, $1/3,\,2/9,\,4/27,\dots$, form a geometric sequence of ratio $2/3$ (the
value of $Q$) whose sum is $1$, and the charged-lepton phase is close to its second term. We fixed in advance a set of
16 candidate phases for the neutrinos (the terms $t_n$, their partial sums, and $2/9+t_n$, $n\le6$). With $Q=2/3$, the
measured ratio $\Delta m^2_{21}/\Delta m^2_{32}=0.0307\pm0.0008$ allows a phase in $[0.4754,0.4825]$ (or its mirror
image); one candidate falls inside, $1/3+4/27=13/27$. With 16 candidates such a hit happens by chance about once in
ten, so this is not significant by itself. Taken exactly, it fixes the ratio to $0.031230$ and, with the measured
$\Delta m^2_{32}$,
\begin{equation}
  \Delta m^2_{21}=(7.667\pm0.087)\times10^{-5}\ \text{eV}^2 .
\end{equation}
After this value was written, we compared it with the first JUNO measurement~\cite{juno},
$(7.50\pm0.12)\times10^{-5}$\,eV$^2$: a difference of $1.1$ combined $\sigma$. JUNO's coming data will decide.

\section{An earlier, rival prediction}
The same author published on 30 October 2026, from a larger network of relations~\cite{grmsite}, $m_1=1.55$\,meV and
$\Sigma m_\nu=60.7$\,meV (normal ordering), with $m_1$ fixed by a different relation (the geometric mean of the Planck
energy and the energy of the final horizon temperature). The two rules cannot hold together: in the network they are in
tension by more than $40\,\sigma$. They agree on the normal ordering and on a sum close to the minimum allowed by
oscillations; the mass of the lightest neutrino will tell them apart.

\section{How these predictions can fail}
\begin{center}\small
\begin{tabular}{lll}
\toprule
quantity & loop & refuted if \\
\midrule
mass ordering & normal & inverted ordering established \\
$\Sigma m_\nu$ & $59.35\pm0.30$\,meV & settles clearly away (${<}55$ or ${>}80$\,meV) \\
$m_1$ & $0.365$\,meV & measured well above $0.4$\,meV \\
$m_\beta$ & $8.84$\,meV & a signal above ${\sim}15$\,meV (today ${<}450$\,meV~\cite{katrin}) \\
$m_{\beta\beta}$ & $1.1$--$4.1$\,meV & a signal clearly above $4.1$\,meV \\
$m_\tau$ & $1776.969$\,MeV & a measurement settles ${>}3\,\sigma$ away \\
$\Delta m^2_{21}$ (Cantor) & $7.667\pm0.087\times10^{-5}$ & ratio outside $0.03123\pm0.0002$ \\
\bottomrule
\end{tabular}
\end{center}
The absence of a $0\nu\beta\beta$ signal would not support the loop, since the neutrino may not be its own antiparticle.

\section{Discussion}
Mathematically, the loop is Koide's relation; its geometric form is Foot's, and its application to neutrinos with one
amplitude in opposition is Brannen's idea. What this note adds is a reading (a balance with no seed and no scale), a
derivation of the neutrino masses from the measured differences alone with no adjusted phase, an explicit statement of
the tension between exact $Q=2/3$ and exact $\delta=2/9$, and predictions archived with a timestamp before the data
that can refute them. The quark masses do not follow the relation ($Q\approx0.85$ and $0.73$ for up- and down-type,
with scheme-dependent masses), and we make no claim about them. Code, tests and the dated files are available
at~\cite{grm}.

\paragraph{Acknowledgements.} The computations, the code and the drafting of this note were done with the assistance of
an AI system (Claude, Anthropic). The ideas of the loop and of the balance between full and void are the author's.

\begin{thebibliography}{9}
\bibitem{koide} Y.~Koide, \emph{A fermion-boson composite model of quarks and leptons}, Phys. Lett. B \textbf{120}
  (1983) 161; \emph{New view of quark and lepton mass hierarchy}, Phys. Rev. D \textbf{28} (1983) 252.
\bibitem{foot} R.~Foot, \emph{A note on Koide's lepton mass relation}, arXiv:hep-ph/9402242 (1994).
\bibitem{brannen} C.~Brannen, \emph{Koide's mass formula for neutrinos}, APS Northwest Section meeting (2006),
  \url{https://meetings-archive.aps.org/nws/2006/b3/14}.
\bibitem{zenczykowski} P.~\.Zenczykowski, \emph{Koide's $Z_3$-symmetric parametrization, quark masses and mixings},
  arXiv:1301.4143 (2013).
\bibitem{pdg} Particle Data Group, Review of Particle Physics (2024), summary tables.
\bibitem{desi} DESI Collaboration, \emph{DESI 2024 VI: Cosmological constraints from the measurements of baryon
  acoustic oscillations}, arXiv:2404.03002 (2024).
\bibitem{juno} JUNO Collaboration, \emph{First measurement of reactor neutrino oscillations at JUNO},
  arXiv:2511.14593 (2025).
\bibitem{katrin} KATRIN Collaboration, \emph{Direct neutrino-mass measurement based on 259 days of KATRIN data},
  Science (2025).
\bibitem{grm} A.~Charretier, \emph{GRM: a Global Relational Model (grm-core)}, Zenodo,
  doi:\href{https://doi.org/10.5281/zenodo.23964101}{10.5281/zenodo.23964101} (2026).
\bibitem{grmsite} A.~Charretier, \emph{Global Relational Model}, \url{https://grm.anthony-charretier.fr} (dated
  predictions of 30 October 2026).
\end{thebibliography}
\end{document}
